\section{Methodology}

This section explains the full workflow and process used to develop this model. The methodology is divided into three main stages: data collection, data preparation, and model training/evaluation. Data was both synthetically generated and collected from archival Kepler sources. The collected data were then processed into fixed-length light-curve samples, labelled by class, and split into training, validation, and testing sets. Finally, a convolutional neural network was trained on the data to distinguish between confirmed exoplanet transit signals and eclipsing binary light curves. The models performance is then evaluated using classification metrics and further misclassification analysis.

To make the workflow reproducible, the project was organized as a staged pipeline. Configuration files were used to store parameters such as data paths, model settings, and evaluation choices, while DVC was used to track the order of processing stages and their outputs.

\subsection{Data Collection}

Two sets of data were used for the research in this paper. First synthetic data was developed with simple equations. This allowed us to appropriately test and check the CNN functions correctly under controlled conditions. By introducing simplistic and harder data sets the model can be evaluated solely on its set up and not by errors in the data. Afterwards, real light curve data is collected from the \textit{Kepler} mission to utilize the CNN for classification. The following subsections explain the generation and collection of all data used in this paper.

\subsubsection{Synthetic Data}

For the purpose of confirming the formed model works without error, synthetic data was made. This was developed in three stages with each being more and more realistic. The first stage contains straight lines with only dips representing transits and eclipsing binary cases. The model is expected to come back with nearly perfect results. The generated data used randomized values for the transit or eclipsing binary dips, transit/eclipsing durations, full orbit periods all within typical ranges expected of the real life systems. Each instance contains 4894 data points which is roughly 100 days of data taking the time step to be \textit{Kepler's} cadence of 29.4 minutes. As well, the shape of dips is different for the two cases with transits having the characteristic U-shape and eclipsing binaries having the V-shape to match real data.

The second stage contains the same parameters as the initial stage but now with added noise. The noise is generated through a uniform random distribution where each data point is given a small white noise value. This creates a more difficult classification task than the first stage and tests whether the model can still identify the injected signals when handling less idealized light curves.

The final synthetic stage is the most realistic and acts as a bridge between the generated signals and observed data. Since the injected signal type is known exactly, the synthetic set is useful for testing whether the model architecture and preprocessing pipeline behave as expected. However, these synthetic data are not used as the main scientific test of the classifier; they are instead used as a control before applying the CNN to real archival light curves. \textbf{(AI)}

\subsubsection{Kepler Extraction}

The real Kepler data collection was designed to produce two labelled classes for supervised learning: confirmed transiting exoplanets and confirmed eclipsing binaries. Confirmed exoplanets could be selected directly from the NASA Exoplanet Archive using the KOI disposition label, while eclipsing binaries required a separate catalogue of known EB KepIDs. The KepID was therefore used as the common identifier linking catalogue labels, archive metadata, and the downloaded light curves. Additional metadata columns were also retained so that model errors could later be compared against physical and observational properties such as signal-to-noise ratio, transit depth, and orbital period. \textbf{(AI)}

Data was collected from the \textbf{(Give better naming?)} catalogue from two tables. The first is \textbf{(give a label with the link in the footnote)} where information on all KOI's are kept. The second is solely for eclipsing binaries and contains all of their KepIDs. KepID's (Kepler Identification \textbf{(Is this the expanded word?)}) is a unique number code for each object of interest. By using the second table, we are able to collect the full table data from the first table. 

\begin{lstlisting}[language=SQL, caption={SQL query used to extract confirmed Kepler exoplanets data from Kepler Exoplanet Archive.}, label={lst:transit-query}]
SELECT {kep_query_cols}
FROM cumulative
WHERE koi_disposition = 'CONFIRMED'
\end{lstlisting}

In order to collect the necessary data, a simple SQL query is used first for the confirmed exoplanets. Listed above in Listing~\ref{lst:transit-query}, this query contains a few labels. First kep\_query\_cols contains the following columns from the Kepler archive listed in Table~\ref{tab:kep_query_cols}. The name of the table is cumulative in the NASA Exoplanet Archive TAP service and so we use this naming. Finally as this is for confirmed planets only we specify the Exoplanet Archive Disposition to be "CONFIRMED". 

\begin{table}[]
    \centering
    \begin{tabular}{|c|c|}\hline
       Kepler Archive Column Name  & SQL Table Name \\\hline
        KepID & kepid \\\hline
        Exoplanet Archive Disposition & koi\_disposition\\\hline
        Transit Signal to Noise & koi\_model\_snr \\\hline
        Transit Depth [ppm] & koi\_depth \\\hline
        Orbital Period [days] & koi\_period
    \\ \hline\end{tabular}
    \caption{This table lists all columns extracted from the Kepler Exoplanet Archive. The left column tells the label NASA provides in their table while the right is the SQL name used in querying.}
    \label{tab:kep_query_cols}
\end{table}

Now for the eclipsing binaries, simply taking the disposition as "Eclipsing Binary" or "False Positive" does not work as the first is not an option and the second includes multiple possible cases. In order to gather the data a second table is required. This is the Kepler Eclipsing Binaries table and it contains the KepID's of all the EBs confirmed from the \textit{Kepler} mission. A simple web-scrape of the website provides the table data and after taking the KepID's another simple SQL query is used written out in Listing~\ref{lst:eb-query}. The list of KepID's is saved into kepid\_string and we simply specify this in the query to pull only the EB's and their column info. 

\begin{lstlisting}[language=SQL, caption={SQL query used to extract eclipsing binary information from Kepler Exoplanet Archive.}, label={lst:eb-query}]
SELECT {kep_query_cols}
FROM cumulative
WHERE kepid IN ({kepid_string})
\end{lstlisting}


\subsection{Data Preparation}

The overall process of downloading each light curve, and saving its data appropriately is done using the \textit{lightkurve} Python package \textbf{(GET CITATION)}. Each KepID is one by one inputted into the library to collect its flux data. As \textbf{Kepler} works in sectors, there are several data sets per star and thus we must stitch them all together. Next all NaN values are removed, outliers with a separation of $5\sigma$ were dropped as well. In the same step the data is normalized to a base flux of 1. Next the time-series data was flattened which is done to remove slow, long-term trends from the light curves while trying to preserve the shorter transit/eclipse events. This is done in the \textit{lightkurve} package by applying a Savitzky-Golay filter. The full \textit{Kepler} objects flux is then saved to a NumPy array.

To restrict errors in the process, the total batch is split into splits of size 200. Every 200 Kepler objects is saved into a csv file and so on until all light curves are processed. Only two columns are left over as we only need one for holding the list of data points as well as saving the associated KepID of each row. As now there are multiple rows per KepID it is important to keep them labelled for the modelling step. 

A total of \textbf{(Make sure)} 2745 confirmed exoplanets are pulled from the Kepler Exoplanet Archive catalogue. This same process is repeated for the eclipsing binaries. The initial file we collect from NASA has 2878 \textbf{(Also check this one)} confirmed eclipsing binaries. Once again the KepID's are split into chunks of 200 and saved per csv file. 

After collection, the raw data and catalogue metadata needed to be converted into consistent formats suitable for the model and future evaluation. First, each input into the CNN model must be the same length and so we needed to determine how to limit the size of each light curve. We also wanted to expand the total number of training instances we could use which would improve the final metrics. The method chosen was to split each light curve into chunks of element length 4894. Given the time step of the \textit{Kepler} mission was 29.4 minutes, this means each chunk covers a length of $\sim 100$ days. Most confirmed planets can be split into three or more of these chunks increasing the total instances to \textbf{(How many?)}. Afterwards, the chunks are saved into new csv's, again with each containing 200 total \textit{Kepler} objects.

The data now consists of two tables, one for the exoplanets and another for the eclipsing binaries. Each table contains the columns listed in the left column of Table~\ref{tab:kep_query_cols}. For the purpose of further analysis, after the model is done learning, we will need the signal to noise ratio of each test instance. While Kepler does provide a column, some rows especially in the eclipsing binary cases were missing values. In order to include SNR for all test cases, a simple formula is applied and is shown in Equation~\ref{eq:snr_calc}. This column is much more consistent in the table and so this allows us to come up with an SNR scale for the analysis of the model results. Using NumPy for the standard deviation and the transit depth values, the results for each instance is added to the updated table. 

\begin{align}
    \text{SNR} = \frac{\text{Transit Depth}}{\text{Standard Deviation}} \label{eq:snr_calc}
\end{align}

For the light curve data we needed to split the data into three parts. A training, validation, and testing set. First we make sure there are no overlaps in the exoplanet and eclipsing binary datasets. While this should not happen it is a good precaution to make sure no data is repeated in both classes. This would damage the efficiency of the model in correctly classifying future light curves. Next, training works best when the data set of trained on exoplanets and eclipsing binaries is the same. This gives equal importance to learning the patterns in each class. Thus the data sets are made to match the size of the smaller one. This takes the eclipsing binary dataset from \textbf{(How many total?)} to \textbf{(How many total?)}. 

Now the data is ready to be split into the different sets. However it is very important to manage how the instances are split. For one KepID, we cannot have one of its chunks in the training set, while another is in the validation or test set. This is unfair to the model as it will likely recognize the exact pattern giving a false boost in the final algorithms score. This would be a form of over fitting which would be notable in the validation results. Thus to avoid this, we split the data into the three sets not by random but by the KepID's.

The data is set to put 80\% of the data into the training set. So in the preparation, the data takes 80\% of the KepID's and saves them to the training set. The remaining 20\% is split in half to give 10\% to validation and 10\% to testing. Instead of taking the first 80\% of KepID's, a shuffle function was developed which shuffles around the data while keeping the KepID's together. This allows us to randomize which KepID's end up being used for training, validation, and testing while maintaining the rule of keeping associated KepID's in the same set. The randomization is set to a seed so the separation is reproducible for the exact results in this paper. 

This overall process is done to the exoplanet and eclipsing binary data sets separately to make the shuffling step more concrete. Before concatenating the two classes, the arrays of the class values needs to be made. Class 0 is the label given to all confirmed exoplanets while for eclipsing binaries they are labelled class 1. Now the training, validation, and test inputs and outputs are combined giving the completed sets of data and their outputs. Each of the three sets now properly separated are given one last shuffle to randomize the order of class 0 and class 1 inputs for the model. Once more, a set seed is used to maintain reproducibility.  

\subsection{Model}

Table~\ref{tab:layer_order} shows the finalized order of layers used in the neural network. The input layer takes in a stacked numpy array with each of the light curves being analysed. From here each light curve is put into a convolutional layer which runs a kernal of size 11 across the entire time series. Rectified Linear Units (ReLU) activations are used here to limit risks of learning decay \textbf{(Is that what is called?)}. The next layers are pooling which takes the largest element in a group of 2 and saves it discarding the other. This lowers the computational load on the algorithm as it learns while retaining the important information. Following this a dropout layer is applied. These remove the usage of random neurons in the layer. This is done to keep the network from relying too heavily on one path or unit. This then reduces over fitting and helps to increase confidence in our results. The second convolution layer now has twice as many channels as the first convolution layer. This means each of the 32 channels from the first layer is split into 64 channel each with an applied filter learning and reading for patterns in the data. This gives a total of 64x32 channels totalling 2048 results to analyze dips and trends in the light curves. After the process is repeated with the max pooling and dropout, a global pooling layer is used to \textbf{(does what actually?)}. Now the results extracted from the 2048 channels is put into a dense layer then one more dropout at a higher rate. Lastly the output layer comes back with a single neuron which gives its guess of the light curve classification. 0 if it is certain the data shows a transiting exoplanet and 1 if certain it represents an eclipsing binary. 

\begin{table}[h]
    \centering
    \begin{tabular}{|c|l|}\hline
           Layer Type&  Parameters\\\hline
           Input&  \\\hline
           Convolution& 32 channels, kernal size = 11, same padding, ReLU activation \\\hline
           Max Pooling& Pool Size = 2 \\\hline
           Dropout& Rate = 0.2 \\\hline
           Convolution&  64 channels, kernal size = 11, same padding, ReLU activation\\\hline
           Max Pooling& Pool Size = 2 \\\hline
           Dropout& Rate = 0.2 \\\hline
           Global Pooling &  \\\hline
           Dense&  \\\hline
           Dropout& Rate = 0.3 \\\hline
           Output&  \\ \hline
    \end{tabular}
    \caption{Order of all layers used in the training of the neural network. All parameters applied are listed as well.}
    \label{tab:layer_order}
\end{table}

\textbf{(Does some of this info go to the results section or is it good here?)}
Upon running the model through with data, the history is saved into a keras file. This can then be referred back to and loaded to observe how the model performed at each epoch in its training process. Several metrics are chosen to record in the learning. Accuracy was chosen to represent how well the model does at prediction the correct class overall. As well, the loss is recorded to show if the model is making decisions with confidence or with blind guesses. Both have results recorded at each epoch for both the training data and validation data. The training results in loss and accuracy are taken into account for updating the model and the validation tells if the model is over or under fitting the data.

Several other metrics are recorded upon the completion of the learning. Once the test data is used against the finalized model, a file is saved holding the recall, precision, F1 score, and ROC AUC. These results can then be monitored to understand more concretely what the model is good at and where it struggles.

To further analyze the CNN, all misclassified light curves are saved in a compressed numpy (.npz) file. Their KepID, input, true class, predicted class, and confidence provided by the model are all saved. This is then split into all false positives and all false negatives to understand better the limitations of the current model. 

\subsubsection{Analysis Setup \textbf{(Def need a better title name)}}

In order to analyze the mistakes of the model, several steps are taken to visyalize the models errors. \textbf{(Do I need to write anything here?)}
